\IfFileExists{stacks-project.cls}{%
\documentclass{stacks-project}
}{%
\documentclass{amsart}
}

% For dealing with references we use the comment environment
\usepackage{verbatim}
\newenvironment{reference}{\comment}{\endcomment}
%\newenvironment{reference}{}{}
\newenvironment{slogan}{\comment}{\endcomment}
\newenvironment{history}{\comment}{\endcomment}

% For commutative diagrams we use Xy-pic
\usepackage[all]{xy}

% We use 2cell for 2-commutative diagrams.
\xyoption{2cell}
\UseAllTwocells

% We use multicol for the list of chapters between chapters
\usepackage{multicol}

% This is generally recommended for better output
\usepackage{lmodern}
\usepackage[T1]{fontenc}

% For cross-file-references
\usepackage{xr-hyper}

% Package for hypertext links:
\usepackage{hyperref}

% For any local file, say "hello.tex" you want to link to please
% use \externaldocument[hello-]{hello}
\externaldocument[introduction-]{introduction}
\externaldocument[conventions-]{conventions}
\externaldocument[sets-]{sets}
\externaldocument[categories-]{categories}
\externaldocument[topology-]{topology}
\externaldocument[sheaves-]{sheaves}
\externaldocument[sites-]{sites}
\externaldocument[stacks-]{stacks}
\externaldocument[fields-]{fields}
\externaldocument[algebra-]{algebra}
\externaldocument[brauer-]{brauer}
\externaldocument[homology-]{homology}
\externaldocument[derived-]{derived}
\externaldocument[simplicial-]{simplicial}
\externaldocument[more-algebra-]{more-algebra}
\externaldocument[smoothing-]{smoothing}
\externaldocument[modules-]{modules}
\externaldocument[sites-modules-]{sites-modules}
\externaldocument[injectives-]{injectives}
\externaldocument[cohomology-]{cohomology}
\externaldocument[sites-cohomology-]{sites-cohomology}
\externaldocument[dga-]{dga}
\externaldocument[dpa-]{dpa}
\externaldocument[sdga-]{sdga}
\externaldocument[hypercovering-]{hypercovering}
\externaldocument[schemes-]{schemes}
\externaldocument[constructions-]{constructions}
\externaldocument[properties-]{properties}
\externaldocument[morphisms-]{morphisms}
\externaldocument[coherent-]{coherent}
\externaldocument[divisors-]{divisors}
\externaldocument[limits-]{limits}
\externaldocument[varieties-]{varieties}
\externaldocument[topologies-]{topologies}
\externaldocument[descent-]{descent}
\externaldocument[perfect-]{perfect}
\externaldocument[more-morphisms-]{more-morphisms}
\externaldocument[flat-]{flat}
\externaldocument[groupoids-]{groupoids}
\externaldocument[more-groupoids-]{more-groupoids}
\externaldocument[etale-]{etale}
\externaldocument[chow-]{chow}
\externaldocument[intersection-]{intersection}
\externaldocument[pic-]{pic}
\externaldocument[weil-]{weil}
\externaldocument[adequate-]{adequate}
\externaldocument[dualizing-]{dualizing}
\externaldocument[duality-]{duality}
\externaldocument[discriminant-]{discriminant}
\externaldocument[derham-]{derham}
\externaldocument[local-cohomology-]{local-cohomology}
\externaldocument[algebraization-]{algebraization}
\externaldocument[curves-]{curves}
\externaldocument[resolve-]{resolve}
\externaldocument[models-]{models}
\externaldocument[functors-]{functors}
\externaldocument[equiv-]{equiv}
\externaldocument[pione-]{pione}
\externaldocument[etale-cohomology-]{etale-cohomology}
\externaldocument[proetale-]{proetale}
\externaldocument[relative-cycles-]{relative-cycles}
\externaldocument[more-etale-]{more-etale}
\externaldocument[trace-]{trace}
\externaldocument[crystalline-]{crystalline}
\externaldocument[spaces-]{spaces}
\externaldocument[spaces-properties-]{spaces-properties}
\externaldocument[spaces-morphisms-]{spaces-morphisms}
\externaldocument[decent-spaces-]{decent-spaces}
\externaldocument[spaces-cohomology-]{spaces-cohomology}
\externaldocument[spaces-limits-]{spaces-limits}
\externaldocument[spaces-divisors-]{spaces-divisors}
\externaldocument[spaces-over-fields-]{spaces-over-fields}
\externaldocument[spaces-topologies-]{spaces-topologies}
\externaldocument[spaces-descent-]{spaces-descent}
\externaldocument[spaces-perfect-]{spaces-perfect}
\externaldocument[spaces-more-morphisms-]{spaces-more-morphisms}
\externaldocument[spaces-flat-]{spaces-flat}
\externaldocument[spaces-groupoids-]{spaces-groupoids}
\externaldocument[spaces-more-groupoids-]{spaces-more-groupoids}
\externaldocument[bootstrap-]{bootstrap}
\externaldocument[spaces-pushouts-]{spaces-pushouts}
\externaldocument[spaces-chow-]{spaces-chow}
\externaldocument[groupoids-quotients-]{groupoids-quotients}
\externaldocument[spaces-more-cohomology-]{spaces-more-cohomology}
\externaldocument[spaces-simplicial-]{spaces-simplicial}
\externaldocument[spaces-duality-]{spaces-duality}
\externaldocument[formal-spaces-]{formal-spaces}
\externaldocument[restricted-]{restricted}
\externaldocument[spaces-resolve-]{spaces-resolve}
\externaldocument[formal-defos-]{formal-defos}
\externaldocument[defos-]{defos}
\externaldocument[cotangent-]{cotangent}
\externaldocument[examples-defos-]{examples-defos}
\externaldocument[algebraic-]{algebraic}
\externaldocument[examples-stacks-]{examples-stacks}
\externaldocument[stacks-sheaves-]{stacks-sheaves}
\externaldocument[criteria-]{criteria}
\externaldocument[artin-]{artin}
\externaldocument[quot-]{quot}
\externaldocument[stacks-properties-]{stacks-properties}
\externaldocument[stacks-morphisms-]{stacks-morphisms}
\externaldocument[stacks-limits-]{stacks-limits}
\externaldocument[stacks-cohomology-]{stacks-cohomology}
\externaldocument[stacks-perfect-]{stacks-perfect}
\externaldocument[stacks-introduction-]{stacks-introduction}
\externaldocument[stacks-more-morphisms-]{stacks-more-morphisms}
\externaldocument[stacks-geometry-]{stacks-geometry}
\externaldocument[moduli-]{moduli}
\externaldocument[moduli-curves-]{moduli-curves}
\externaldocument[examples-]{examples}
\externaldocument[exercises-]{exercises}
\externaldocument[guide-]{guide}
\externaldocument[desirables-]{desirables}
\externaldocument[coding-]{coding}
\externaldocument[obsolete-]{obsolete}
\externaldocument[fdl-]{fdl}
\externaldocument[index-]{index}

% Theorem environments.
%
\theoremstyle{plain}
\newtheorem{theorem}[subsection]{Theorem}
\newtheorem{proposition}[subsection]{Proposition}
\newtheorem{lemma}[subsection]{Lemma}

\theoremstyle{definition}
\newtheorem{definition}[subsection]{Definition}
\newtheorem{example}[subsection]{Example}
\newtheorem{exercise}[subsection]{Exercise}
\newtheorem{situation}[subsection]{Situation}

\theoremstyle{remark}
\newtheorem{remark}[subsection]{Remark}
\newtheorem{remarks}[subsection]{Remarks}

\numberwithin{equation}{subsection}

% Macros
%
\def\lim{\mathop{\mathrm{lim}}\nolimits}
\def\colim{\mathop{\mathrm{colim}}\nolimits}
\def\Spec{\mathop{\mathrm{Spec}}}
\def\Hom{\mathop{\mathrm{Hom}}\nolimits}
\def\Ext{\mathop{\mathrm{Ext}}\nolimits}
\def\SheafHom{\mathop{\mathcal{H}\!\mathit{om}}\nolimits}
\def\SheafExt{\mathop{\mathcal{E}\!\mathit{xt}}\nolimits}
\def\Sch{\mathit{Sch}}
\def\Mor{\mathop{\mathrm{Mor}}\nolimits}
\def\Ob{\mathop{\mathrm{Ob}}\nolimits}
\def\Sh{\mathop{\mathit{Sh}}\nolimits}
\def\NL{\mathop{N\!L}\nolimits}
\def\CH{\mathop{\mathrm{CH}}\nolimits}
\def\proetale{{pro\text{-}\acute{e}tale}}
\def\etale{{\acute{e}tale}}
\def\QCoh{\mathit{QCoh}}
\def\Ker{\mathop{\mathrm{Ker}}}
\def\Im{\mathop{\mathrm{Im}}}
\def\Coker{\mathop{\mathrm{Coker}}}
\def\Coim{\mathop{\mathrm{Coim}}}

% Boxtimes
%
\DeclareMathSymbol{\boxtimes}{\mathbin}{AMSa}{"02}

%
% Macros for moduli stacks/spaces
%
\def\QCohstack{\mathcal{QC}\!\mathit{oh}}
\def\Cohstack{\mathcal{C}\!\mathit{oh}}
\def\Spacesstack{\mathcal{S}\!\mathit{paces}}
\def\Quotfunctor{\mathrm{Quot}}
\def\Hilbfunctor{\mathrm{Hilb}}
\def\Curvesstack{\mathcal{C}\!\mathit{urves}}
\def\Polarizedstack{\mathcal{P}\!\mathit{olarized}}
\def\Complexesstack{\mathcal{C}\!\mathit{omplexes}}
% \Pic is the operator that assigns to X its picard group, usage \Pic(X)
% \Picardstack_{X/B} denotes the Picard stack of X over B
% \Picardfunctor_{X/B} denotes the Picard functor of X over B
\def\Pic{\mathop{\mathrm{Pic}}\nolimits}
\def\Picardstack{\mathcal{P}\!\mathit{ic}}
\def\Picardfunctor{\mathrm{Pic}}
\def\Deformationcategory{\mathcal{D}\!\mathit{ef}}
\usepackage{cleveref}
\crefname{originallemma}{original lemma}{original lemmas}
\Crefname{originallemma}{Original lemma}{Original lemmas}
\makeatletter
\newlabel{lemma-adequate-local}{{06VI}{source}{}{sourcepointer-1}{}}
\newlabel{lemma-adequate-local@cref}{{[originallemma][0][]06VI}{[1][0][]source}{}{}{}}
\makeatother
\crefname{originallemma}{original lemma}{original lemmas}
\Crefname{originallemma}{Original lemma}{Original lemmas}
\makeatletter
\newlabel{lemma-cokernel-adequate}{{06V2}{source}{}{sourcepointer-2}{}}
\newlabel{lemma-cokernel-adequate@cref}{{[originallemma][0][]06V2}{[1][0][]source}{}{}{}}
\makeatother
\crefname{originallemma}{original lemma}{original lemmas}
\Crefname{originallemma}{Original lemma}{Original lemmas}
\makeatletter
\newlabel{lemma-colimit-adequate}{{06V4}{source}{}{sourcepointer-3}{}}
\newlabel{lemma-colimit-adequate@cref}{{[originallemma][0][]06V4}{[1][0][]source}{}{}{}}
\makeatother
\crefname{originallemma}{original lemma}{original lemmas}
\Crefname{originallemma}{Original lemma}{Original lemmas}
\makeatletter
\newlabel{lemma-extension-adequate}{{06V8}{source}{}{sourcepointer-4}{}}
\newlabel{lemma-extension-adequate@cref}{{[originallemma][0][]06V8}{[1][0][]source}{}{}{}}
\makeatother
\crefname{originallemma}{original lemma}{original lemmas}
\Crefname{originallemma}{Original lemma}{Original lemmas}
\makeatletter
\newlabel{lemma-kernel-adequate}{{06V3}{source}{}{sourcepointer-5}{}}
\newlabel{lemma-kernel-adequate@cref}{{[originallemma][0][]06V3}{[1][0][]source}{}{}{}}
\makeatother
\crefname{originallemma}{original lemma}{original lemmas}
\Crefname{originallemma}{Original lemma}{Original lemmas}
\makeatletter
\newlabel{lemma-same-cohomology-adequate}{{06VM}{source}{}{sourcepointer-6}{}}
\newlabel{lemma-same-cohomology-adequate@cref}{{[originallemma][0][]06VM}{[1][0][]source}{}{}{}}
\makeatother
\crefname{originallemma}{original lemma}{original lemmas}
\Crefname{originallemma}{Original lemma}{Original lemmas}
\makeatletter
\newlabel{lemma-sheafification-adequate}{{06VP}{source}{}{sourcepointer-7}{}}
\newlabel{lemma-sheafification-adequate@cref}{{[originallemma][0][]06VP}{[1][0][]source}{}{}{}}
\makeatother
\title{The category of adequate modules is abelian and its inclusion is exact}
\author{The Stacks Project Authors}
\date{stacks-project:adequate.tex:local-snapshot-20261012; GNU FDL 1.2 or later}
\begin{document}
\maketitle
\noindent\textit{Editorial note.} All delivered mathematical text is the author's own, copied from the pinned Stacks Project chapter file adequate.tex; no statement or proof step was rewritten, reordered or removed.\par\medskip
\noindent\textit{Editorial note.} The counted spans carry navigation stubs for the references leaving them; no bibliography citation occurs in the spans.\par\medskip
\noindent\textit{Editorial note.} The preamble span points at the chapter's own preamble.tex shipped with the frozen source, used verbatim as staged.\par\medskip
\noindent\textit{Allowed standard prerequisites (editorial).}\begin{itemize}
\item Kernels and cokernels of maps of adequate functors are adequate. Same chapter, Lemmas lemma-kernel-adequate and lemma-cokernel-adequate, invoked at proof lines 2053-2056. Give adequacy of the kernel and cokernel in the abelian category of functors.
\item Adequacy is preserved by sheafification, and the associated functor construction. Same chapter, Lemma lemma-sheafification-adequate, invoked at proof lines 2060-2064. Transfers adequacy from the presheaf cokernel to the cokernel.
\item Vanishing of H\textasciicircum{}1 for adequate modules on affine schemes, and exactness of the associated functor sequence. Same chapter, Lemma lemma-same-cohomology-adequate, invoked at proof lines 2075-2077. Makes the sequence of functors associated to a short exact sequence exact.
\item Adequacy is local, and an extension of adequate functors is adequate. Same chapter, Lemmas lemma-adequate-local and lemma-extension-adequate, invoked at proof lines 2078-2081. Gives item (3).
\item Colimits of adequate functors and their adequacy. Same chapter, Lemma lemma-colimit-adequate, invoked at proof lines 2084-2099. Gives item (4).
\end{itemize}\medskip
\ifdefined\hypertarget\hypertarget{sourcepointer-1}{}\fi
\noindent\textit{Source reference.} Original lemma 06VI: sources/stacks-adequate/pinned-original-source.tex, source lines 1707--1707. Pointer outside the counted span, invoked once for context; the load-bearing argument of the counted statement is the author's proof reproduced in full in the proof span.\par\medskip
\ifdefined\hypertarget\hypertarget{sourcepointer-2}{}\fi
\noindent\textit{Source reference.} Original lemma 06V2: sources/stacks-adequate/pinned-original-source.tex, source lines 684--684. Pointer outside the counted span, invoked once for context; the load-bearing argument of the counted statement is the author's proof reproduced in full in the proof span.\par\medskip
\ifdefined\hypertarget\hypertarget{sourcepointer-3}{}\fi
\noindent\textit{Source reference.} Original lemma 06V4: sources/stacks-adequate/pinned-original-source.tex, source lines 734--734. Pointer outside the counted span, invoked once for context; the load-bearing argument of the counted statement is the author's proof reproduced in full in the proof span.\par\medskip
\ifdefined\hypertarget\hypertarget{sourcepointer-4}{}\fi
\noindent\textit{Source reference.} Original lemma 06V8: sources/stacks-adequate/pinned-original-source.tex, source lines 910--910. Pointer outside the counted span, invoked once for context; the load-bearing argument of the counted statement is the author's proof reproduced in full in the proof span.\par\medskip
\ifdefined\hypertarget\hypertarget{sourcepointer-5}{}\fi
\noindent\textit{Source reference.} Original lemma 06V3: sources/stacks-adequate/pinned-original-source.tex, source lines 709--709. Pointer outside the counted span, invoked once for context; the load-bearing argument of the counted statement is the author's proof reproduced in full in the proof span.\par\medskip
\ifdefined\hypertarget\hypertarget{sourcepointer-6}{}\fi
\noindent\textit{Source reference.} Original lemma 06VM: sources/stacks-adequate/pinned-original-source.tex, source lines 1896--1896. Pointer outside the counted span, invoked once for context; the load-bearing argument of the counted statement is the author's proof reproduced in full in the proof span.\par\medskip
\ifdefined\hypertarget\hypertarget{sourcepointer-7}{}\fi
\noindent\textit{Source reference.} Original lemma 06VP: sources/stacks-adequate/pinned-original-source.tex, source lines 1985--1985. Pointer outside the counted span, invoked once for context; the load-bearing argument of the counted statement is the author's proof reproduced in full in the proof span.\par\medskip
\begin{lemma}
\label{lemma-abelian-adequate}
Let $S$ be a scheme.
\begin{enumerate}
\item The category $\textit{Adeq}(\mathcal{O})$ is abelian.
\item The functor
$\textit{Adeq}(\mathcal{O}) \to
\textit{Mod}((\Sch/S)_\tau, \mathcal{O})$
is exact.
\item If $0 \to \mathcal{F}_1 \to \mathcal{F}_2 \to \mathcal{F}_3 \to 0$
is a short exact sequence of $\mathcal{O}$-modules and
$\mathcal{F}_1$ and $\mathcal{F}_3$ are adequate, then
$\mathcal{F}_2$ is adequate.
\item The category $\textit{Adeq}(\mathcal{O})$ has colimits and
$\textit{Adeq}(\mathcal{O}) \to
\textit{Mod}((\Sch/S)_\tau, \mathcal{O})$
commutes with them.
\end{enumerate}
\end{lemma}

\begin{proof}
Let $\varphi : \mathcal{F} \to \mathcal{G}$ be a map of adequate
$\mathcal{O}$-modules. To prove (1) and (2) it suffices to show that
$\mathcal{K} = \Ker(\varphi)$ and
$\mathcal{Q} = \Coker(\varphi)$ computed in
$\textit{Mod}((\Sch/S)_\tau, \mathcal{O})$ are adequate.
Let $U = \Spec(A)$ be an affine scheme over $S$.
Let $F = F_{\mathcal{F}, A}$ and $G = F_{\mathcal{G}, A}$.
By
Lemmas \ref{lemma-kernel-adequate} and
\ref{lemma-cokernel-adequate}
the kernel $K$ and cokernel $Q$ of the induced map
$F \to G$ are adequate functors.
Because the kernel is computed on the level of presheaves, we see
that $K = F_{\mathcal{K}, A}$ and we conclude $\mathcal{K}$ is adequate.
To prove the result for the cokernel, denote $\mathcal{Q}'$ the presheaf
cokernel of $\varphi$. Then $Q = F_{\mathcal{Q}', A}$ and
$\mathcal{Q} = (\mathcal{Q}')^\#$. Hence $\mathcal{Q}$
is adequate by
Lemma \ref{lemma-sheafification-adequate}.

\medskip\noindent
Let $0 \to \mathcal{F}_1 \to \mathcal{F}_2 \to \mathcal{F}_3 \to 0$
is a short exact sequence of $\mathcal{O}$-modules and
$\mathcal{F}_1$ and $\mathcal{F}_3$ are adequate.
Let $U = \Spec(A)$ be an affine scheme over $S$.
Let $F_i = F_{\mathcal{F}_i, A}$. The sequence of functors
$$
0 \to F_1 \to F_2 \to F_3 \to 0
$$
is exact, because for $V = \Spec(B)$ affine over $U$ we have
$H^1(V, \mathcal{F}_1) = 0$ by
Lemma \ref{lemma-same-cohomology-adequate}.
Since $F_1$ and $F_3$ are adequate functors by
Lemma \ref{lemma-adequate-local}
we see that $F_2$ is adequate by
Lemma \ref{lemma-extension-adequate}.
Thus $\mathcal{F}_2$ is adequate.

\medskip\noindent
Let $\mathcal{I} \to \textit{Adeq}(\mathcal{O})$, $i \mapsto \mathcal{F}_i$
be a diagram. Denote $\mathcal{F} = \colim_i \mathcal{F}_i$
the colimit computed in
$\textit{Mod}((\Sch/S)_\tau, \mathcal{O})$.
To prove (4) it suffices to show that $\mathcal{F}$ is adequate.
Let $\mathcal{F}' = \colim_i \mathcal{F}_i$ be the colimit computed
in presheaves of $\mathcal{O}$-modules. Then
$\mathcal{F} = (\mathcal{F}')^\#$.
Let $U = \Spec(A)$ be an affine scheme over $S$.
Let $F_i = F_{\mathcal{F}_i, A}$. By
Lemma \ref{lemma-colimit-adequate}
the functor $\colim_i F_i = F_{\mathcal{F}', A}$ is adequate.
Lemma \ref{lemma-sheafification-adequate}
shows that $\mathcal{F}$ is adequate.
\end{proof}

\end{document}
